\documentclass[11pt,a4paper]{amsart}

\usepackage[T1]{fontenc}
\usepackage{amsmath,amssymb,amsthm,mathtools}
\usepackage{enumitem}
\usepackage{microtype}
\usepackage{etoolbox}
\usepackage[margin=2.7cm]{geometry}
\usepackage{xcolor}
\definecolor{linkcol}{rgb}{0.10,0.20,0.50}
\usepackage[colorlinks=true,linkcolor=linkcol,citecolor=linkcol,urlcolor=linkcol]{hyperref}
\makeatletter
\g@addto@macro\UrlBreaks{\do\_\do\.}
\makeatother

\theoremstyle{plain}
\newtheorem{theorem}{Theorem}[section]
\newtheorem{proposition}[theorem]{Proposition}
\newtheorem{lemma}[theorem]{Lemma}
\newtheorem{corollary}[theorem]{Corollary}
\theoremstyle{definition}
\newtheorem{definition}[theorem]{Definition}
\newtheorem{example}[theorem]{Example}
\newtheorem{remark}[theorem]{Remark}

% Keep consecutive theorem/proof blocks visually distinct in the compact amsart layout.
\AfterEndEnvironment{theorem}{\par\smallskip}
\AfterEndEnvironment{proposition}{\par\smallskip}
\AfterEndEnvironment{lemma}{\par\smallskip}
\AfterEndEnvironment{corollary}{\par\smallskip}
\AfterEndEnvironment{proof}{\par\smallskip}

\newcommand{\RR}{\mathbb R}
\newcommand{\CC}{\mathbb C}
\newcommand{\NN}{\mathbb N}
\newcommand{\one}{\mathbf 1}
\newcommand{\rank}{\operatorname{rank}}
\newcommand{\corank}{\operatorname{corank}}
\newcommand{\indm}{\operatorname{ind}_{-}}
\newcommand{\Spec}{\operatorname{Spec}}
\newcommand{\diag}{\operatorname{diag}}
\newcommand{\spanr}{\operatorname{span}_{\RR}}
\newcommand{\Res}{\mathcal R}
\newcommand{\Mcal}{\mathcal M}
\newcommand{\Hcal}{\mathcal H}
\newcommand{\Tcal}{\mathcal T}
\newcommand{\Ecal}{\mathcal E}
\emergencystretch=1.5em

\begin{document}

\title[The Stokes--Cartan one-defect law]{At most Lorentzian:\
a universal Stokes--Cartan signature law in genus zero}

\author{M.~Haufschild}
\address{Independent researcher}
\email{haufschildmichael@gmail.com}
\date{3 August 2026}

\subjclass[2020]{Primary 34M40, 14D24; Secondary 15A18, 32G34, 16G20}
\keywords{irregular Riemann--Hilbert correspondence, wild nonabelian Hodge theory,
Stokes data, fission tree, generalized Cartan matrix, Lorentzian signature, effective resistance,
principal conormal, Rees degeneration}

\begin{abstract}
Genus-zero wild nonabelian Hodge spaces carry symmetric Cartan forms assembled from local Stokes
contacts and ordinary type-$A$ legs.  We prove that every such form has negative index at most one.
More precisely, for ramification vector $r$, a rooted fission-contact tree with positive edge
lengths $\ell_e$, descendant clusters $A_e$, and an auxiliary root height $H>1$, the core Cartan
matrix admits the canonical factorization
\[
 C=I+\sum_e\ell_e u_eu_e^T-(H-1)rr^T,
 \qquad u_e=(r_i\one_{A_e}(i))_i.
\]
Thus every local fission channel is a positive principal-conormal square and total ramified rank is
the unique possible negative channel.  If
\[
 P=I+\sum_e\ell_eu_eu_e^T,\qquad
 c=r^TP^{-1}r,\qquad \Omega=c^{-1}-(H-1),
\]
then $\Omega$ is independent of the artificial root extension and completely determines inertia:
$C$ is positive definite, affine of corank one, or Lorentzian according as $\Omega$ is positive,
zero, or negative.  We also prove $\det C=\det(P)c\Omega$.

The reciprocal capacity $c^{-1}$ is the effective resistance of the contact tree.  Series and
parallel reduction gives an exact corona-pruning recursion, and turns root independence into the
cancellation of two equal series increments.  Attaching arbitrary finite type-$A$ legs changes
the positive diagonal from $1$ to $1/(l_i+1)$ and cannot create another defect.  Consequently every
full type-$A$ genus-zero diagram obeys the same trichotomy.  For an adjacency diagram $B$ this is
the spectral obstruction $\lambda_2(B)\leq2$; hence higher-signature Kac--Moody diagrams do not
occur in this class.  Finally, the wild-character-variety dimension becomes one global rank-square
budget minus explicit positive leg and conormal energies.  The result is a diagram-level law: it
does not identify the conormal coefficients with characteristic-cycle or vanishing-cycle
multiplicities.
\end{abstract}

\maketitle

\section{Introduction}

Wild nonabelian Hodge theory relates meromorphic connections, irregular Higgs bundles, and Stokes
local systems on curves \cite{BiquardBoalch2004}.  In genus zero, the formal data determine a finite
diagram with a symmetric generalized Cartan matrix.  This matrix controls the quadratic dimension
formula for the associated wild character variety \cite{BoalchYamakawa2020,Doucot2021Diagrams}.
The diagram is invariant under the symplectic automorphisms of the Weyl algebra, including
Fourier--Laplace transformation \cite{Doucot2021Diagrams}.  Fission trees record how the formal
structure group breaks at an irregular pole and classify admissible deformation data
\cite{Boalch2008Kac,BoalchDoucotRembado2025,Boalch2025Counting}.

The resulting graphs range beyond finite and affine Dynkin diagrams.  Hyperbolic examples are
known, and the question of which Kac--Moody root systems can arise is part of the current
classification problem \cite{Boalch2012Simply,Boalch2025Counting}.  Ramification makes that problem
sharper: core diagrams may have loops or signed entries, while some twisted irregular classes
nevertheless yield honest graphs that are not fission graphs \cite{Doucot2025Graphs}.  Ultrametric
inequalities constrain their rescaled edge multiplicities, but those inequalities alone do not
state the inertia of the Cartan form.

This paper proves a universal one-defect law.  Its shortest form is
\begin{equation}\label{eq:intro-law}
 \boxed{\quad C=P-\lambda rr^T,\qquad P>0.\quad}
\end{equation}
The positive part $P$ is not an unspecified matrix.  It is the sum of a positive diagonal and one
rank-one Gram term for every fission cut.  The only subtractive term is the square of total
ramified rank.  Formula \eqref{eq:intro-law} immediately suggests negative index at most one, but
its force is greater: a single intrinsic scalar $\Omega$ detects all three possible signatures,
the affine kernel, the determinant, and the exact transition wall.

For the core matrix $B=(B_{ij})$ and ramification orders $r_i$, introduce the complementary
contact coordinates
\begin{equation}\label{eq:intro-K}
 K_{ij}=1+\frac{B_{ij}}{r_ir_j}\quad(i\ne j),\qquad
 K_{ii}=1+\frac{B_{ii}-1}{r_i^2}.
\end{equation}
The fission-contact hierarchy realizes $K_{ij}$ as the height of the meet of leaves $i$ and $j$.
Extend the root to any height $H>\max(1,K_{ij})$.  If $A_e$ is the descendant leaf set of an edge,
$\ell_e>0$ its contact length, and $u_e=(r_i\one_{A_e}(i))_i$, our first main theorem is
\begin{equation}\label{eq:intro-factor}
 C=I+\sum_e\ell_eu_eu_e^T-(H-1)rr^T.
\end{equation}
The extension to height $H$ is auxiliary.  Adding a root segment of length $t$ adds $t rr^T$ to
the positive Gram part and simultaneously replaces $H-1$ by $H-1+t$, leaving $C$ unchanged.

Set
\begin{equation}\label{eq:intro-capacity}
 P=I+\sum_e\ell_eu_eu_e^T,\qquad
 c=r^TP^{-1}r,\qquad
 \Omega=c^{-1}-(H-1).
\end{equation}
Then $c>0$, and every $x$ splits uniquely into a vector of zero total ramified rank and a multiple
of $P^{-1}r$.  In these coordinates,
\begin{equation}\label{eq:intro-centered}
 x^TCx=z^TPz+\frac{(r^Tx)^2}{c^2}\,c^2\Omega.
\end{equation}
Consequently $C$ is positive definite for $\Omega>0$, positive semidefinite with a one-dimensional
kernel for $\Omega=0$, and nonsingular with exactly one negative eigenvalue for $\Omega<0$.
Neither multiple affine directions nor coexistence of affine and Lorentz directions is possible.
Throughout, ``Lorentzian'' means signature $(N-1,1)$ for a nonsingular real symmetric form on an
$N$-dimensional space; no stronger condition on all proper principal submatrices is built into the
term.

There are four reasons the factorization is geometric rather than merely a rank-one perturbation
observation.
\begin{enumerate}[label=(\roman*),leftmargin=2.2em]
\item Each positive summand is the square of the principal conormal
$\nu_e(x)=\sum_{i\in A_e}r_ix_i$ of a canonical cluster hyperplane.
\item The reciprocal capacity $c^{-1}$ is the effective resistance of the rooted contact tree,
with leaf resistance $r_i^{-2}$.  Corona pruning is therefore exact series--parallel reduction.
\item Ordinary type-$A$ legs admit an exact path completion.  Their Schur complement replaces the
unit leaf resistance by $((l_i+1)r_i^2)^{-1}$ and preserves the one-defect law.
\item The global dimension formula becomes an exact conservation law: one total-rank gain minus
positive leg fluctuations and positive fission-conormal losses.
\end{enumerate}

For honest graphs the theorem has a concise spectral consequence.  If $B$ is the adjacency matrix
and its eigenvalues are ordered decreasingly, then
\[
                         \lambda_2(B)\le2.
\]
If the spectral radius is below two the Cartan form is finite; if it equals two the eigenvalue two
is simple; if it exceeds two there is exactly one eigenvalue above two and none equal to two.  Thus
the known hyperbolic examples are the maximal indefinite possibility: higher-signature
Kac--Moody systems are excluded from type-$A$ genus-zero nonabelian Hodge diagrams.

The matrix-theory ancestors of the argument are the weighted-tree descriptions and inverse
$M$-matrix properties of ultrametric and generalized ultrametric matrices
\cite{MartinezMichonSanMartin1994,McDonaldEtAl1995,NabbenVarga1995}.  A recent treatment of shifted
generalized ultrametric matrices and tree-structured inverses is \cite{Nabben2026Acyclic}.  Those
works concern matrix recognition and inverses.  The new content here is the exact conversion from
Stokes contact to the Cartan factorization \eqref{eq:intro-factor}, the intrinsic capacity wall,
the full-leg and global genus-zero passage, the conormal/dimension law, and the spectral exclusion
for nonabelian Hodge diagrams.  We state the distinction because a tree Gram decomposition by
itself is classical and would not constitute the claimed result.

The paper is organized as follows.  Section~\ref{sec:geometry} isolates the precise geometric
input.  Section~\ref{sec:factor} proves the one-defect factorization.  Section~\ref{sec:capacity}
proves the capacity trichotomy and root invariance.  Section~\ref{sec:corona} gives the electrical
corona recursion.  Section~\ref{sec:legs} treats arbitrary type-$A$ legs.  Sections
\ref{sec:conormal} and \ref{sec:dimension} give conormal/Rees and dimension formulations.
Section~\ref{sec:global} passes to global genus-zero diagrams, and
Section~\ref{sec:spectral} derives graph-theoretic consequences and examples.  The final section
records exact scope and falsifiers.

\section{Stokes contact and rooted fission data}\label{sec:geometry}

We separate the geometric input from the linear algebra.  This prevents a pairwise ultrametric
inequality from being silently upgraded to a simultaneous tree realization.

\subsection{Decorated cores}

Let $I$ be a nonempty finite set.  A \emph{decorated core} consists of a symmetric matrix
$B=(B_{ij})_{i,j\in I}$ and positive integers $r_i$.  In the Stokes application, $I$ indexes
Stokes circles, $r_i$ is the ramification order, and $B$ is the core matrix of the nonabelian
Hodge diagram.  Put
\begin{equation}\label{eq:Cartan-def}
                         C=2I-B.
\end{equation}
Diagonal entries of $B$ are retained; thus this notation includes generalized diagrams with
loops.  Define $K$ by \eqref{eq:intro-K}.  Both cases combine into
\begin{equation}\label{eq:B-contact}
            B_{ij}=\delta_{ij}+r_ir_j(K_{ij}-1).
\end{equation}

\begin{definition}\label{def:rooted-realization}
A \emph{rooted contact realization} of $K$ is a finite rooted tree $\Tcal$ whose leaves are
labelled by $I$, together with a real height $\kappa(v)$ at every retained vertex, strictly
increasing along each edge from a leaf toward the root, such that
\[
                         K_{ij}=\kappa(i\wedge j)
\]
for all $i,j\in I$.  Here $i\wedge i=i$, so leaf heights encode the diagonal contacts.
\end{definition}

For irregular classes, the realization is the common-truncation skeleton of the fission tree.
The edge and loop formulas for a Stokes circle give \eqref{eq:B-contact}; common Puiseux
truncations give the meet rule.  The following elementary reconstruction records why one common
tree, not a collection of pairwise witnesses, is present.

\begin{lemma}[Contact dendrogram]\label{lem:dendrogram}
Suppose a symmetric matrix $K$ satisfies $K_{ii}<K_{ij}$ for $i\ne j$ and, for every three
distinct indices, the two larger of $K_{ij},K_{ik},K_{jk}$ are equal.  Then $K$ has a rooted
contact realization, unique after unary vertices are suppressed.  Its clusters are the classes of
the relations
\[
 i\sim_a j\quad\Longleftrightarrow\quad i=j\ \text{or}\ K_{ij}\le a
\]
as $a$ ranges over the off-diagonal values of $K$.
\end{lemma}

\begin{proof}
Fix $a$.  Symmetry and reflexivity are clear.  If $i\sim_a j$ and $j\sim_a k$, then the stated
three-point condition forbids $K_{ik}$ from being strictly larger than both $K_{ij}$ and $K_{jk}$;
hence $K_{ik}\le a$.  Thus $\sim_a$ is an equivalence relation.  Its classes are nested as $a$
increases.  Make every non-singleton class a vertex, attach it to the least strictly larger class
that contains it, and attach the labelled singleton leaves at height $K_{ii}$.  Give a class born
at level $a$ height $a$.  The strict diagonal inequality makes pendant edges positive; deleting
repeated levels suppresses unary vertices.  The smallest cluster containing $i$ and $j$ is born at
$K_{ij}$, so its vertex is their meet.  Conversely, any rooted realization has exactly these
threshold clusters, proving uniqueness.
\end{proof}

\begin{proposition}[Geometric input]\label{prop:geometric-input}
The complementary-contact matrix of a finite type-$A$ irregular class has a rooted contact
realization.  For a genus-zero nonabelian Hodge diagram, any generic one-pole reading supplies such
a realization of its core, and the remaining vertices form ordinary type-$A$ legs.
\end{proposition}

\begin{proof}
At one irregular pole, the common truncations of the exponential factors form the fission tree.
Choose a level $\eta$ above all retained characteristic exponents.  On a common truncation branch,
let $R(u)>0$ be the cumulative ramification index at height $u$.  Telescoping the explicit loop and
edge formulae gives the complementary coarea coordinate
\begin{equation}\label{eq:contact-coarea}
 \kappa(v)=\eta-\int_{h(v)}^\eta\frac{du}{R(u)}.
\end{equation}
Profiles on two branches agree above their meet, so \eqref{eq:contact-coarea} is independent of
the leaf used to read a retained common-truncation vertex.  Formula
\eqref{eq:contact-coarea} is the summation-by-parts form of the rescaled source formula
\cite[Lemma~4.3]{Doucot2025Graphs}; the diagonal case is its loop analogue.  The diagonal formula reads
$K_{ii}=\kappa(i)$, and the pair formula reads $K_{ij}=\kappa(i\wedge j)$.  If $v$ is strictly below
$w$ on one branch, then
\[
 \kappa(w)-\kappa(v)=\int_{h(v)}^{h(w)}\frac{du}{R(u)}>0.
\]
Thus contact is strictly increasing toward the root.  This gives
Definition~\ref{def:rooted-realization}; equivalently one may apply Lemma~\ref{lem:dendrogram} to
the rescaled ultrametric relations proved from the explicit Stokes edge formulae in
\cite[Theorem~4.2]{Doucot2025Graphs}.

For a general algebraic connection on the affine line, the diagram is invariant under the
symplectic Weyl-algebra action \cite{Doucot2021Diagrams}.  The generic class reconstructed from the
enriched tree has all Stokes circles at infinity, while the enriched tree and diagram are invariant
under the symplectic action \cite[Theorems~1.4--1.6]{Doucot2026Basic}.  This supplies the required
one-pole reading.  The conjugacy-class data attach the ordinary type-$A$ legs used in the diagram
construction \cite{BoalchYamakawa2020,Doucot2026Basic}.  Only equality of the diagram is used here;
no claim that all weak representations have already been proved analytically isomorphic is needed.
\end{proof}

\begin{remark}
Proposition~\ref{prop:geometric-input} is the geometric boundary of the paper.  The remaining
proofs use only a rooted realization and the exact identity \eqref{eq:B-contact}.  Conversely, a
symmetric matrix satisfying an isolated collection of inequalities is not declared realizable by
an irregular connection.
\end{remark}

\subsection{Root extension and cut coordinates}

Choose $H>\max(1,\kappa(v))$ and attach a new root of height $H$.  Orient edges away from that
root.  For each edge $e$, let $A_e\subseteq I$ be its descendant leaves and let $\ell_e>0$ be the
drop in height along $e$.  Define
\begin{equation}\label{eq:cut-vectors}
 u_e=(r_i\one_{A_e}(i))_{i\in I},\qquad r=(r_i)_{i\in I},
 \qquad \lambda=H-1.
\end{equation}
The root extension is deliberately noncanonical.  Section~\ref{sec:capacity} will show that every
quantity claimed to be intrinsic is independent of it.

\section{The Stokes--Cartan one-defect factorization}\label{sec:factor}

\begin{lemma}[Root-path Gram identity]\label{lem:root-path}
For all leaves $i,j$,
\begin{equation}\label{eq:root-path}
 H-K_{ij}=\sum_e\ell_e\one_{A_e}(i)\one_{A_e}(j).
\end{equation}
Consequently the matrix $(H-K_{ij})$ is positive semidefinite.
\end{lemma}

\begin{proof}
An edge belongs to both root-to-leaf paths precisely when both leaves are descendants of that
edge.  The intersection of the two paths ends at $i\wedge j$.  Telescoping the positive height
drops along this common path gives $H-\kappa(i\wedge j)=H-K_{ij}$, proving
\eqref{eq:root-path}.  The right side is the Gram matrix
$\sum_e\ell_e a_ea_e^T$ with $a_e=(\one_{A_e}(i))_i$.
\end{proof}

\begin{theorem}[Universal one-defect law]\label{thm:one-defect}
For every rooted decorated Stokes core,
\begin{equation}\label{eq:one-defect-matrix}
 \boxed{\quad C=I+\sum_e\ell_eu_eu_e^T-\lambda rr^T.\quad}
\end{equation}
Equivalently, for every $x\in\RR^I$,
\begin{equation}\label{eq:one-defect-q}
 \boxed{\quad
 x^TCx=\sum_i x_i^2+
 \sum_e\ell_e\left(\sum_{i\in A_e}r_ix_i\right)^2
 -(H-1)\left(\sum_i r_ix_i\right)^2.\quad}
\end{equation}
\end{theorem}

\begin{proof}
By \eqref{eq:B-contact}, \eqref{eq:Cartan-def}, and Lemma~\ref{lem:root-path},
\begin{align*}
 C_{ij}
 &=\delta_{ij}+r_ir_j(1-K_{ij})\\
 &=\delta_{ij}+r_ir_j(H-K_{ij})-(H-1)r_ir_j\\
 &=\delta_{ij}+\sum_e\ell_e
   \bigl(r_i\one_{A_e}(i)\bigr)\bigl(r_j\one_{A_e}(j)\bigr)
   -\lambda r_ir_j.
\end{align*}
This is \eqref{eq:one-defect-matrix} entrywise.  Multiplying by $x_ix_j$ and summing gives
\eqref{eq:one-defect-q}.
\end{proof}

For later use, write
\begin{equation}\label{eq:mass-conormals}
 M(x)=r^Tx=\sum_i r_ix_i,
 \qquad \nu_e(x)=u_e^Tx=\sum_{i\in A_e}r_ix_i.
\end{equation}

\begin{corollary}[Canonical positive hyperplane]\label{cor:hyperplane}
The Cartan form is positive definite on
\[
                         \Hcal_r=\ker M.
\]
In particular $\indm(C)\le1$.
\end{corollary}

\begin{proof}
On $\Hcal_r$, the final term of \eqref{eq:one-defect-q} vanishes, while
$\sum_i x_i^2>0$ for $x\ne0$.  A quadratic form positive on a codimension-one subspace cannot have
a negative subspace of dimension two.
\end{proof}

The corollary is already the qualitative signature law.  The next section identifies the exact
wall and rules out a simultaneous zero and negative direction.

\section{Capacity, determinant, and the trichotomy}\label{sec:capacity}

Put
\begin{equation}\label{eq:P-c-Omega}
 P=I+\sum_e\ell_eu_eu_e^T,\qquad
 c=r^TP^{-1}r,\qquad
 \Omega=c^{-1}-\lambda.
\end{equation}
The matrix $P$ is positive definite, so $c>0$.

\begin{theorem}[Capacity normal form]\label{thm:capacity-normal}
For $x\in\RR^I$, set $m=M(x)$ and
\begin{equation}\label{eq:centered-z}
 z=x-\frac{m}{c}P^{-1}r.
\end{equation}
Then $r^Tz=0$, the decomposition is unique, and
\begin{equation}\label{eq:capacity-completion}
 \boxed{\quad x^TCx=z^TPz+m^2\Omega.\quad}
\end{equation}
\end{theorem}

\begin{proof}
Since $r^TP^{-1}r=c$, equation \eqref{eq:centered-z} gives $r^Tz=0$.  If
$x=z+aP^{-1}r$ with $r^Tz=0$, applying $r^T$ gives $a=m/c$, proving uniqueness.  Moreover
$z^TP(P^{-1}r)=z^Tr=0$.  Hence
\begin{align*}
 x^TCx
 &=x^TPx-\lambda(r^Tx)^2\\
 &=z^TPz+\frac{m^2}{c^2}r^TP^{-1}r-\lambda m^2\\
 &=z^TPz+m^2(c^{-1}-\lambda).
\end{align*}
\end{proof}

\begin{theorem}[Finite--affine--Lorentz trichotomy]\label{thm:trichotomy}
Exactly one of the following holds.
\begin{enumerate}[label=(\roman*),leftmargin=2.2em]
\item If $\Omega>0$, then $C$ is positive definite.
\item If $\Omega=0$, then $C$ is positive semidefinite,
\[
             \ker C=\spanr(P^{-1}r),\qquad \corank C=1.
\]
\item If $\Omega<0$, then $C$ is nonsingular and $\indm(C)=1$.
\end{enumerate}
In particular, $C$ has at most one nonpositive eigenvalue.  A negative and a zero eigenvalue can
never coexist.
\end{theorem}

\begin{proof}
In \eqref{eq:capacity-completion}, $z^TPz\ge0$, with equality exactly when $z=0$.
If $\Omega>0$, equality of the whole expression forces both $z=0$ and $m=0$, hence $x=0$.
If $\Omega=0$, the zero vectors are exactly those with $z=0$, namely the stated line.
If $\Omega<0$, the vector $P^{-1}r$ is negative because its centered part is zero and its mass is
$c>0$.  Corollary~\ref{cor:hyperplane} bounds the negative index by one, so it equals one.  A kernel
vector would have to be independent from this negative line or lie on it; the first possibility
would create a nonpositive plane and the second is excluded by strict negativity.  Equivalently,
the determinant formula below is nonzero.
\end{proof}

\begin{theorem}[Determinant wall]\label{thm:determinant}
One has
\begin{equation}\label{eq:determinant}
                   \boxed{\quad\det C=\det(P)c\Omega.\quad}
\end{equation}
Thus the affine wall is exactly $\Omega=0$, and it is crossed transversely in the scalar defect
$\lambda$.
\end{theorem}

\begin{proof}
The matrix determinant lemma gives
\[
 \det(P-\lambda rr^T)=\det(P)(1-\lambda r^TP^{-1}r)
 =\det(P)(1-\lambda c)=\det(P)c(c^{-1}-\lambda).
\]
\end{proof}

\begin{proposition}[Root-extension gauge invariance]\label{prop:root-invariance}
Extend the auxiliary root by length $t>0$.  Then
\begin{equation}\label{eq:root-update}
 P_t=P+t rr^T,\qquad \lambda_t=\lambda+t,\qquad
 c_t=\frac{c}{1+tc}.
\end{equation}
Consequently
\begin{equation}\label{eq:root-invariants}
 c_t^{-1}-\lambda_t=\Omega,\qquad
 \det(P_t)c_t=\det(P)c.
\end{equation}
In particular the trichotomy and determinant formula are intrinsic to the decorated diagram.
\end{proposition}

\begin{proof}
The new edge has descendant set $I$, so its cut vector is $r$, proving the first two updates.
Sherman--Morrison gives
\[
 P_t^{-1}=P^{-1}-\frac{tP^{-1}rr^TP^{-1}}{1+tc},
\]
and contraction with $r$ gives the capacity update.  Therefore
$c_t^{-1}=c^{-1}+t$, proving the first invariant.  The determinant lemma gives
$\det(P_t)=\det(P)(1+tc)$, which proves the second.
\end{proof}

\begin{remark}[Intrinsic formulation]
Choose any allowed root extension and compute $\Omega$.  Proposition~\ref{prop:root-invariance}
shows the result is independent of the choice.  Thus $\Omega$ is not a height normalization; it is
a scalar invariant of the Stokes--Cartan datum.  It is the signed distance between effective
resistance and global defect, as the next section makes literal.
\end{remark}

\section{Corona pruning and effective resistance}\label{sec:corona}

The term \emph{corona} will mean the terminal branches adjacent to a chosen internal vertex.  To
prune a corona is to eliminate all variables inside it while retaining the total current through
its stem.  The elimination is an exact Schur complement, not an asymptotic approximation.

Change variables from $x_i$ to leaf currents
\begin{equation}\label{eq:current-change}
                         m_i=r_ix_i.
\end{equation}
Then the positive core energy becomes
\begin{equation}\label{eq:electrical-energy}
 x^TPx=\sum_i\frac{m_i^2}{r_i^2}
       +\sum_e\ell_e\left(\sum_{i\in A_e}m_i\right)^2.
\end{equation}
Interpret $r_i^{-2}$ as a terminal resistance and $\ell_e$ as an edge resistance.  The current in
edge $e$ is the total leaf current below it.

\begin{lemma}[Parallel completion]\label{lem:parallel}
Let $R_1,\dots,R_s>0$ and $m_1+\cdots+m_s=M$.  Put
\[
 S=\sum_{a=1}^sR_a^{-1},\qquad R_{\parallel}=S^{-1}.
\]
Then
\begin{equation}\label{eq:parallel-completion}
 \sum_{a=1}^sR_am_a^2
 =\sum_{a=1}^sR_a\left(m_a-\frac{R_a^{-1}}{S}M\right)^2
  +R_{\parallel}M^2.
\end{equation}
\end{lemma}

\begin{proof}
Expand the squares.  The cross terms sum to $-2M^2/S$ and the constant terms to $M^2/S$.
Rearrangement gives \eqref{eq:parallel-completion}.  Positivity also shows that the unique
energy-minimizing split is proportional to conductance.
\end{proof}

\begin{definition}[Recursive resistance]\label{def:recursive-R}
Assign a leaf $i$ the resistance $R_i=r_i^{-2}$.  If an internal vertex $v$ has children $w$, and
the edge from $v$ to $w$ has length $\ell_{vw}$, define
\begin{equation}\label{eq:resistance-recursion}
                  R_v=\left(\sum_{w\prec v}
                    \frac{1}{\ell_{vw}+R_w}\right)^{-1}.
\end{equation}
The value at the chosen root is $R_{\mathrm{eff}}$.
\end{definition}

\begin{theorem}[Corona-capacity law]\label{thm:corona-capacity}
Successive pruning by \eqref{eq:resistance-recursion} gives
\begin{equation}\label{eq:min-energy}
 \min_{\sum_i m_i=M}x^TPx=R_{\mathrm{eff}}M^2,
 \qquad R_{\mathrm{eff}}=c^{-1}.
\end{equation}
Moreover
\begin{equation}\label{eq:Omega-resistance}
                         \Omega=R_{\mathrm{eff}}-\lambda.
\end{equation}
A root extension adds the same series resistance $t$ to both terms in
\eqref{eq:Omega-resistance}.  Thus root invariance is the series-cancellation law
\[
 (R_{\mathrm{eff}}+t)-(\lambda+t)=R_{\mathrm{eff}}-\lambda.
\]
\end{theorem}

\begin{proof}
At a terminal corona, the resistance from a child subtree is its effective resistance plus the
stem resistance $\ell_{vw}$ in series.  Lemma~\ref{lem:parallel} replaces all child branches by
their parallel resistance and a sum of positive fluctuation squares.  Induction toward the root
proves the first minimization formula.

On the other hand, Theorem~\ref{thm:capacity-normal} applied to the positive form alone gives
\[
 x^TPx=z^TPz+M^2/c
\]
under the constraint $r^Tx=M$.  Its unique minimum is $M^2/c$.  Comparing the two minima gives
$R_{\mathrm{eff}}=c^{-1}$, and \eqref{eq:Omega-resistance} follows from definition.
The final assertion is ordinary series addition together with
Proposition~\ref{prop:root-invariance}.
\end{proof}

\begin{corollary}[Local-to-global computation]\label{cor:message-passing}
The signature phase of a rooted core is determined by one bottom-up pass through the fission tree:
compute $R_v$ by \eqref{eq:resistance-recursion} and compare $R_{\mathrm{eff}}$ with $H-1$.
Every pruned corona contributes only positive fluctuation squares; no local pruning step can create
a second negative direction.
\end{corollary}

\begin{remark}[Vanishing-cycle language]
Formula \eqref{eq:parallel-completion} is an algebraic elimination law for a terminal cluster.  It
is reasonable to view the discarded positive fluctuation coordinates as quadratic
``vanishing-cycle channels.''  We do not identify them with sheaf-theoretic vanishing cycles, nor
do we claim that $\ell_e$ is a characteristic-cycle multiplicity.  Such an identification would
require an additional microlocal comparison theorem not proved here.
\end{remark}

\section{Type-A legs preserve the one-defect law}\label{sec:legs}

Attach to core vertex $i$ an ordinary path leg of length $l_i\ge0$.  Write its variables as
$y_{i,1},\dots,y_{i,l_i}$, and set
\[
 y_{i,0}=x_i,\qquad y_{i,l_i+1}=0,\qquad
 d_{i,k}=y_{i,k}-y_{i,k+1}\quad(0\le k\le l_i).
\]
Then $\sum_{k=0}^{l_i}d_{i,k}=x_i$.

\begin{lemma}[Path completion]\label{lem:path-completion}
For $d_0+\cdots+d_l=x$,
\begin{equation}\label{eq:path-completion}
 \sum_{k=0}^l d_k^2
 =\sum_{k=0}^l\left(d_k-\frac{x}{l+1}\right)^2+\frac{x^2}{l+1}.
\end{equation}
\end{lemma}

\begin{proof}
Expand the right side and use $\sum d_k=x$.
\end{proof}

Let $\widehat C$ be the full Cartan matrix.  The ordinary type-$A$ edge coupling gives
\begin{equation}\label{eq:full-before-completion}
 \widehat q(x,y)=
 \sum_i\sum_{k=0}^{l_i}d_{i,k}^2
 +\sum_e\ell_e\nu_e(x)^2-\lambda M(x)^2.
\end{equation}
Indeed, the path energy $\sum_kd_{i,k}^2$ equals the leg Cartan energy together with the core
square $x_i^2$ that occurs in \eqref{eq:one-defect-q}.

\begin{theorem}[Full-leg conservation]\label{thm:full-leg}
For arbitrary finite type-$A$ legs,
\begin{align}\label{eq:full-completed}
 \widehat q(x,y)
 &=\sum_i\sum_{k=0}^{l_i}
     \left(d_{i,k}-\frac{x_i}{l_i+1}\right)^2
   +\sum_i\frac{x_i^2}{l_i+1}\notag\\
 &\quad+\sum_e\ell_e\nu_e(x)^2-\lambda M(x)^2.
\end{align}
Consequently $\widehat C$ is positive definite, positive semidefinite of corank one, or has
exactly one negative eigenvalue.  In particular
\begin{equation}\label{eq:full-index}
                 \indm(\widehat C)\le1,\qquad
                 \dim\ker\widehat C\le1.
\end{equation}
\end{theorem}

\begin{proof}
Apply Lemma~\ref{lem:path-completion} to every leg in
\eqref{eq:full-before-completion}.  The fluctuation squares form a positive definite complement
to the variables $x$.  The remaining Schur form is
\begin{equation}\label{eq:leg-Schur}
 S=D+\sum_e\ell_eu_eu_e^T-\lambda rr^T,\qquad
 D=\diag\left(\frac1{l_i+1}\right)>0.
\end{equation}
Replace $I$ by $D$ in Sections~\ref{sec:factor}--\ref{sec:capacity}.  The positive matrix
$P_L=D+\sum_e\ell_eu_eu_e^T$ is still positive definite, so the same centered decomposition and
trichotomy apply.  Completing the legs is a congruence by a positive block and therefore preserves
the nonpositive inertia of $S$.
\end{proof}

\begin{corollary}[Leg capacity]\label{cor:leg-capacity}
For the full diagram, use terminal resistances
\begin{equation}\label{eq:leg-resistance}
                         R_i=\frac{1}{(l_i+1)r_i^2}.
\end{equation}
The recursion \eqref{eq:resistance-recursion} computes the full capacity wall.  In particular,
lengthening a leg changes only a positive terminal resistance and never increases the maximal
negative index.
\end{corollary}

\begin{proof}
Under $m_i=r_ix_i$, the diagonal term in \eqref{eq:leg-Schur} is
$m_i^2/((l_i+1)r_i^2)$.  Apply Theorem~\ref{thm:corona-capacity}.
\end{proof}

\section{Principal conormals and Rees stability}\label{sec:conormal}

The cut form has a direct microlocal interpretation that is exact at the level of linear
conormals.  For an edge $e$, define the cluster hyperplane
\begin{equation}\label{eq:cluster-hyperplane}
             \Hcal_e=\{x\in\RR^I:\nu_e(x)=0\}.
\end{equation}
Its conormal line is generated by the constant covector
\begin{equation}\label{eq:principal-conormal}
                         d\nu_e=u_e.
\end{equation}
The family $\{A_e\}$ is laminar: two descendant sets are disjoint or one contains the other.

\begin{theorem}[Principal-conormal Hessian law]\label{thm:conormal-Hessian}
Let $q_+(x)=x^TPx$.  Its Hessian satisfies
\begin{equation}\label{eq:Hessian-law}
 \frac12\operatorname{Hess}(q_+)
 =I+\sum_e\ell_e\,d\nu_e\otimes d\nu_e.
\end{equation}
For the full diagram, $I$ is replaced by
$D=\diag((l_i+1)^{-1})$.  Thus the complete local Stokes contribution is the weighted Gram norm of
the principal conormals of the laminar fission cuts.  The sole negative Hessian channel is
$-\lambda\,dM\otimes dM$.
\end{theorem}

\begin{proof}
This is the polarization of \eqref{eq:one-defect-q} and
\eqref{eq:full-completed}; $d\nu_e=u_e$ and $dM=r$.
\end{proof}

We next record the exact Rees-type stability that follows from this decomposition.  Assign an
indeterminate $s_e$ to each fission edge and define the multiscale family
\begin{equation}\label{eq:Rees-family}
 q_{\boldsymbol s}(x)=x^TDx+
   \sum_e s_e\ell_e\nu_e(x)^2-\lambda M(x)^2.
\end{equation}
At $s_e=1$ this is the reduced full Cartan form.  Setting selected $s_e$ to zero contracts their
graded conormal channels; scaling $s_e=t^{w_e}$ records a weighted contact filtration.  We call
\eqref{eq:Rees-family} the \emph{quadratic multi-Rees family} of the fission filtration.  This
terminology refers only to the explicit filtered quadratic family, not to a claimed equivalence of
categories of Stokes-filtered sheaves.

\begin{corollary}[Rees one-defect stability]\label{cor:Rees-stability}
For every specialization $s_e\ge0$, the form $q_{\boldsymbol s}$ has negative index at most one and
nullity at most one.  It is positive on $\ker M$.  Its phase is determined by
\begin{equation}\label{eq:Rees-gap}
 \Omega(\boldsymbol s)=
 \left[r^T\left(D+\sum_e s_e\ell_eu_eu_e^T\right)^{-1}r\right]^{-1}
 -\lambda.
\end{equation}
\end{corollary}

\begin{proof}
The specialized positive part remains positive definite because $D>0$ and $s_e\ell_e\ge0$.
The proofs of Corollary~\ref{cor:hyperplane} and Theorem~\ref{thm:trichotomy} apply verbatim.
\end{proof}

\begin{remark}
Equation \eqref{eq:Hessian-law} is a principal-conormal statement in a finite linear model.  It
does not identify $\ell_e$ with a multiplicity in the characteristic cycle of an irregular
holonomic $\mathcal D$-module.  Establishing such a statement would require compatibility between
the irregular Riemann--Hilbert correspondence, microlocal specialization, and the fission-contact
filtration.  The present theorem supplies a rigid numerical target for that future comparison.
\end{remark}

\section{The wild dimension conservation law}\label{sec:dimension}

Let $\widehat d=(x,y)$ be a dimension vector for the full diagram, written in the coordinates of
Section~\ref{sec:legs}.  The diagram dimension formula is
\begin{equation}\label{eq:dimension-known}
                \dim\Mcal_B=2-\widehat d^T\widehat C\widehat d
\end{equation}
for the corresponding wild character variety, under the usual nonempty/stability hypotheses
\cite{Doucot2021Diagrams}.

\begin{theorem}[Rank--conormal dimension conservation]\label{thm:dimension-law}
With $d_{i,k}=y_{i,k}-y_{i,k+1}$ as above,
\begin{align}\label{eq:dimension-law}
 \dim\Mcal_B
 &=2+\lambda M(x)^2
 -\sum_i\sum_{k=0}^{l_i}
     \left(d_{i,k}-\frac{x_i}{l_i+1}\right)^2\notag\\
 &\quad-\sum_i\frac{x_i^2}{l_i+1}
 -\sum_e\ell_e\nu_e(x)^2.
\end{align}
This equality is independent of the auxiliary root extension.  It expresses dimension as one
global ramified-rank square budget minus positive leg fluctuations, positive terminal costs, and
positive principal-conormal defects.
\end{theorem}

\begin{proof}
Substitute \eqref{eq:full-completed} into \eqref{eq:dimension-known}.  If the root is extended by
$t$, the new root edge has $A_e=I$ and hence contributes $tM(x)^2$ to the final positive sum, while
$\lambda M(x)^2$ increases by the same amount.  The two changes cancel.
\end{proof}

\begin{corollary}[Relaxed dimension obstruction]\label{cor:dimension-obstruction}
For fixed total ramified rank $M$ and fixed root gauge, the real relaxation of the dimension
functional is maximized by the unique conductance-proportional current split at every corona and by
constant increments along every type-$A$ leg.  Every departure from these balances has an explicit
positive square cost.  An integral dimension vector attains this relaxed bound only when the
balanced profile is integral.
\end{corollary}

\begin{proof}
Apply Lemmas~\ref{lem:parallel} and \ref{lem:path-completion} to
\eqref{eq:dimension-law}.  Equality in either completion gives the stated unique balance.
\end{proof}

This corollary is a useful distinction between the universal law and a sign estimate.  The theorem
does not merely say that certain terms are nonnegative; it locates every loss channel and gives its
equality condition.

\section{Global genus-zero and irregular Riemann--Hilbert invariance}\label{sec:global}

We now formulate the global result.  A \emph{type-$A$ genus-zero nonabelian Hodge diagram} means a
diagram arising from algebraic connection data on a Zariski-open subset of $\mathbb P^1$ through
the standard genus-zero diagram construction, with its ordinary type-$A$ conjugacy-class legs.
This includes ramified and multi-pole formal types covered by the general diagram construction of
\cite{Doucot2021Diagrams}.

\begin{theorem}[Global Stokes--Cartan signature law]\label{thm:global}
Let $\widehat C$ be the symmetric Cartan matrix of a type-$A$ genus-zero nonabelian Hodge diagram.
Then
\begin{equation}\label{eq:global-index}
        \indm(\widehat C)\le1,
 \qquad \dim\ker\widehat C\le1,
 \qquad \indm(\widehat C)+\dim\ker\widehat C\le1.
\end{equation}
Equivalently, every such Cartan form is finite, affine, or Lorentzian.  Its phase, affine null line,
determinant wall, and dimension are given by Theorems
\ref{thm:trichotomy}, \ref{thm:determinant}, and \ref{thm:dimension-law} in any generic one-pole
reading.
\end{theorem}

\begin{proof}
The diagram invariance theorem permits a generic symplectic reading of the same core at infinity
\cite{Doucot2021Diagrams}.  Proposition~\ref{prop:geometric-input} supplies the rooted fission-contact
realization.  Theorem~\ref{thm:one-defect} applies to the core and
Theorem~\ref{thm:full-leg} to all attached type-$A$ legs.  The resulting matrix is the original
diagram Cartan matrix because the diagram, not an asserted analytic isomorphism between all
representations, is the invariant being transported.  The remaining statements follow from the
capacity and dimension theorems.
\end{proof}

\begin{corollary}[Fourier invariance of the wall]\label{cor:Fourier-wall}
The sign of $\Omega$ and the finite/affine/Lorentz phase are invariant under Fourier--Laplace and
the symplectic Weyl-algebra transformations preserving the genus-zero diagram.
\end{corollary}

\begin{proof}
The sign is the inertia of the invariant Cartan matrix by Theorem~\ref{thm:trichotomy}.
\end{proof}

\begin{remark}[Irregular Riemann--Hilbert scope]
Through the irregular Riemann--Hilbert correspondence, the same wild character variety is described
by generalized monodromy and Stokes data on the Betti side.  Theorem~\ref{thm:global} therefore
constrains the Cartan shadow common to these descriptions.  It is not a new proof of the irregular
Riemann--Hilbert equivalence and does not assert that the full Stokes local system can be recovered
from the scalar capacity $\Omega$.
\end{remark}

\section{Graph spectra and Kac--Moody exclusion}\label{sec:spectral}

Suppose the full diagram is an honest graph: $B$ is a real symmetric adjacency matrix with zero
diagonal and nonnegative off-diagonal entries.  Let
$\lambda_1(B)\ge\cdots\ge\lambda_n(B)$ be its eigenvalues.

\begin{theorem}[Second-eigenvalue obstruction]\label{thm:lambda2}
If $n\ge2$ and $B$ is a type-$A$ genus-zero nonabelian Hodge graph, then
\begin{equation}\label{eq:lambda2}
                         \boxed{\lambda_2(B)\le2.}
\end{equation}
More precisely:
\begin{enumerate}[label=(\roman*),leftmargin=2.2em]
\item if $\lambda_1(B)<2$, then $2I-B$ is positive definite;
\item if $\lambda_1(B)=2$, then $2$ is a simple eigenvalue and $2I-B$ is positive semidefinite;
\item if $\lambda_1(B)>2$, then exactly one eigenvalue of $B$ exceeds $2$ and $2$ is not an
eigenvalue.
\end{enumerate}
\end{theorem}

\begin{proof}
The eigenvalues of $2I-B$ are $2-\lambda_i(B)$.  Apply
Theorem~\ref{thm:global}: its negative index is the number of adjacency eigenvalues above two, and
its nullity is the multiplicity of the eigenvalue two.  Formula \eqref{eq:global-index} gives all
claims.
\end{proof}

\begin{corollary}[Kac--Moody signature bound]\label{cor:Kac}
No symmetric generalized Cartan matrix of negative index at least two can be realized by a
type-$A$ genus-zero nonabelian Hodge graph.  Hyperbolic/Lorentzian type is the maximal indefinite
type in this class.
\end{corollary}

The statement is deliberately restricted to the class in Theorem~\ref{thm:global}.  It does not
exclude higher-signature generalized Cartan matrices from other quiver varieties, other genera,
or constructions not admitting the rooted type-$A$ fission reading.

\begin{example}[A sharp twisted graph]\label{ex:twisted-triangle}
Dou\c cot constructs a nonabelian Hodge graph from a twisted irregular class whose three edge
multiplicities are $6,3,4$ \cite[Proposition~3.3]{Doucot2025Graphs}.  In a compatible ordering,
\[
 B=\begin{pmatrix}0&6&4\\6&0&3\\4&3&0\end{pmatrix},\qquad
 C=\begin{pmatrix}2&-6&-4\\-6&2&-3\\-4&-3&2\end{pmatrix}.
\]
For $\one=(1,1,1)^T$,
\[
                         \one^TC\one=6-2(6+4+3)=-20.
\]
Also $\det C=-258$ and $\operatorname{tr}C=6$, independently confirming that the inertia is
$(2,1,0)$.
Thus the universal theorem cannot be strengthened to positivity.  Its conclusion is sharp:
$C$ has exactly one negative eigenvalue and no kernel.  This graph is also useful because it is a
nonabelian Hodge graph that is not an untwisted fission graph, so the phenomenon is not confined to
complete multipartite cores.
\end{example}

\begin{example}[Affine wall]
Whenever a diagram lies on $R_{\mathrm{eff}}=H-1$, its null vector is explicitly
$P^{-1}r$.  In electrical variables, this is the unique conductance-balanced unit-current profile.
The simplicity of the affine eigenvalue is therefore not an appeal to a Dynkin classification; it
is forced directly by strict convexity of the positive Stokes energy.
\end{example}

\begin{example}[Why a second defect is impossible]
Assume two independent vectors had nonpositive Cartan energy.  A nontrivial linear combination of
them has zero total ramified rank.  On that hyperplane the energy is the strictly positive sum of
terminal and cut squares in \eqref{eq:one-defect-q}, a contradiction.  This two-line argument is
the geometric core of the exclusion, while capacity supplies its sharp boundary behavior.
\end{example}

\section{Sharpness, falsifiers, and limits}\label{sec:scope}

We finish by listing what would disprove or extend the theorem.  This is part of the mathematical
statement: a universal-law claim without explicit falsifiers is not auditable.

\subsection{Statements ruled out}

The stronger conjecture ``all nonabelian Hodge Cartan forms are positive semidefinite'' is false by
Example~\ref{ex:twisted-triangle}.  The correct boundary is negative index one.  The theorem also
rules out a weaker but previously plausible phenomenon: a Lorentz direction together with an
independent affine direction.  Capacity completion shows that this cannot occur.

\subsection{Direct falsifiers}

Any one of the following would invalidate the stated global scope.
\begin{enumerate}[label=(\arabic*),leftmargin=2.2em]
\item A realized type-$A$ genus-zero diagram whose Cartan matrix has two negative eigenvalues.
\item A realized diagram whose complementary contacts fail to admit one simultaneous rooted
realization, even though its pairwise inequalities hold.
\item A standard genus-zero diagram containing a leg not obtainable by the ordinary one-ended
type-$A$ attachment used in Section~\ref{sec:legs}.
\item Failure of diagram preservation under the generic symplectic reading used in
Theorem~\ref{thm:global}.
\end{enumerate}
No such example appears in the primary diagram and fission-tree sources cited above.  Conversely,
the theorem is not an analytic existence criterion: a matrix satisfying the one-defect law need
not come from an irregular connection.

\subsection{Boundaries of interpretation}

The conormals in Theorem~\ref{thm:conormal-Hessian} are conormals of explicit cluster hyperplanes in
dimension-vector space.  The corona reduction is an exact Schur complement.  The Rees family is a
filtered deformation of the quadratic form.  None of these statements, alone, computes
microsupports, Stokes sheaf categories, characteristic-cycle coefficients, or vanishing-cycle
multiplicities.  Those are natural next questions because the present law predicts that any
compatible categorification must preserve a single global negative channel while decomposing all
local fission channels positively.

\subsection{Novelty boundary}

Weighted-tree Gram formulae, ultrametric matrices, inverse $M$-matrices, the Stokes diagram, and its
dimension formula are prior results.  The claim of this paper is their previously unstated
intersection: the Stokes normalization \eqref{eq:intro-K} changes $2I-B$ into a positive laminar
Gram form minus exactly one total-rank square; capacity then propagates this identity through
corona pruning, type-$A$ legs, global diagram readings, and the Kac--Moody spectrum.  Removing the
Stokes normalization or the one-pole rooted realization destroys that conclusion.

\section{Conclusion}

The Cartan geometry of a type-$A$ genus-zero wild nonabelian Hodge diagram is at most Lorentzian.
The reason is structural: local irregularity separates along a laminar fission tree and contributes
positive principal-conormal squares, whereas compact genus-zero coupling contributes one square of
total ramified rank with the opposite sign.  Effective resistance compresses the entire tree to a
single intrinsic wall $\Omega$; type-$A$ legs and corona elimination preserve the defect count.

This law supplies a global restriction on the Kac--Moody systems visible in wild nonabelian Hodge
theory and a quantitative decomposition of the dimension formula.  It also isolates a concrete
microlocal problem: construct a categorical irregular Riemann--Hilbert/Rees comparison whose
associated quadratic shadow is the principal-conormal decomposition
\eqref{eq:Hessian-law}.  The theorem proved here is the numerical constraint that any such
comparison must satisfy.

\begin{thebibliography}{99}

\bibitem{BiquardBoalch2004}
O.~Biquard and P.~Boalch,
\emph{Wild non-abelian Hodge theory on curves},
Compos. Math. \textbf{140} (2004), no.~1, 179--204.
\href{https://doi.org/10.1112/S0010437X03000010}{doi:10.1112/S0010437X03000010}.

\bibitem{Boalch2008Kac}
P.~Boalch,
\emph{Irregular connections and Kac--Moody root systems},
arXiv:0806.1050 (2008).
\href{https://arxiv.org/abs/0806.1050}{arXiv:0806.1050}.

\bibitem{Boalch2012Simply}
P.~Boalch,
\emph{Simply-laced isomonodromy systems},
Publ. Math. Inst. Hautes \'{E}tudes Sci. \textbf{116} (2012), 1--68.
\href{https://doi.org/10.1007/s10240-012-0040-2}{doi:10.1007/s10240-012-0040-2}.

\bibitem{Boalch2025Counting}
P.~Boalch,
\emph{Counting the fission trees and nonabelian Hodge graphs (untwisted case)},
J. Geom. Phys. \textbf{214} (2025), 105512.
\href{https://doi.org/10.1016/j.geomphys.2025.105512}{doi:10.1016/j.geomphys.2025.105512}.

\bibitem{BoalchDoucotRembado2025}
P.~Boalch, J.~Dou\c cot and G.~Rembado,
\emph{Twisted local wild mapping class groups: configuration spaces, fission trees and complex
braids},
Publ. Res. Inst. Math. Sci. \textbf{61} (2025), no.~3, 391--445.
\href{https://doi.org/10.4171/PRIMS/61-3-4}{doi:10.4171/PRIMS/61-3-4}.

\bibitem{BoalchYamakawa2020}
P.~Boalch and D.~Yamakawa,
\emph{Diagrams for nonabelian Hodge spaces on the affine line},
C. R. Math. \textbf{358} (2020), no.~1, 59--65.
\href{https://doi.org/10.5802/crmath.11}{doi:10.5802/crmath.11}.

\bibitem{Doucot2021Diagrams}
J.~Dou\c cot,
\emph{Diagrams and irregular connections on the Riemann sphere},
arXiv:2107.02516 (2021).
\href{https://arxiv.org/abs/2107.02516}{arXiv:2107.02516}.

\bibitem{Doucot2025Graphs}
J.~Dou\c cot,
\emph{New nonabelian Hodge graphs from twisted irregular connections},
arXiv:2509.24861 (2025), revised 2026.
\href{https://arxiv.org/abs/2509.24861}{arXiv:2509.24861}.

\bibitem{Doucot2026Basic}
J.~Dou\c cot,
\emph{Basic representations of genus zero nonabelian Hodge spaces},
SIGMA \textbf{22} (2026), 027, 43 pp.
\href{https://doi.org/10.3842/SIGMA.2026.027}{doi:10.3842/SIGMA.2026.027}.

\bibitem{MartinezMichonSanMartin1994}
S.~Mart\'{i}nez, G.~Michon and J.~San Mart\'{i}n,
\emph{Inverse of strictly ultrametric matrices are of Stieltjes type},
SIAM J. Matrix Anal. Appl. \textbf{15} (1994), no.~1, 98--106.
\href{https://doi.org/10.1137/S0895479891217011}{doi:10.1137/S0895479891217011}.

\bibitem{McDonaldEtAl1995}
J.~J.~McDonald, M.~Neumann, H.~Schneider and M.~J.~Tsatsomeros,
\emph{Inverse $M$-matrix inequalities and generalized ultrametric matrices},
Linear Algebra Appl. \textbf{220} (1995), 321--341.
\href{https://doi.org/10.1016/0024-3795(94)00077-Q}{doi:10.1016/0024-3795(94)00077-Q}.

\bibitem{NabbenVarga1995}
R.~Nabben and R.~S.~Varga,
\emph{Generalized ultrametric matrices---a class of inverse $M$-matrices},
Linear Algebra Appl. \textbf{220} (1995), 365--390.
\href{https://doi.org/10.1016/0024-3795(94)00086-S}{doi:10.1016/0024-3795(94)00086-S}.

\bibitem{Nabben2026Acyclic}
R.~Nabben,
\emph{Inverse acyclic $Z$-matrices, matrices of tree structure and shifted generalized
ultrametric matrices},
arXiv:2607.27221 (2026).
\href{https://arxiv.org/abs/2607.27221}{arXiv:2607.27221}.

\end{thebibliography}

\end{document}
